\subsection{عمل المُستقبِل وحدود مدخلاته}
\label{sec:receiver-operation}

يحوّل المُستقبِل مواضع الرد وموضوعات التعليقات الظاهرة إلى الرسالة التي ضمّنها المُرسِل. ويتألف عمله من جزأين: بناء التفسير نفسه للسياق العام، ثم استخراج الاختيارات التي تحملها التعليقات. يجمع هذا القسم خطوات الاستقبال في موضع واحد، بينما يحدد القسم~\ref{sec:multiframe} حقول النقل في نمط التعليقات المتعددة.

\subsubsection{المشاهدات العامة والإعداد المشترك والبيانات الخاصة}
تشمل المدخلات الظاهرة المنشور وتعليقاته ونصوص التعليقات الحاملة، مع معرّفاتها ومعرّفات التعليقات التي ترد عليها. ويتفق الطرفان كذلك على قواعد التحضير: النماذج و\EN{prompts}، واختيار المصادر، والحد الزمني للتعليقات، وترتيب القوائم، وتمثيل الضغط، وإعدادات النقل والحماية. ويتطلب نمط الحماية المعتمد على مفتاح سرًا مشتركًا أيضًا. وتحدد هذه الاتفاقات كيفية تحويل الاختيار الظاهر نفسه إلى بتات.

تختلف قائمة موضوعات التعليق عن قاموس الضغط. تتيح الأولى البدائل التي يستخدمها المُستقبِل لتحديد موضوع التعليق، بينما يفسر الثاني مراجع النص المضغوط بعد استخراج البتات. ويجب أن يطابق كل منهما البنية المقابلة لدى المُرسِل. يستطيع المُستقبِل بناءهما باتباع خطوات التحضير المتفق عليها، أو تلقي بنيتين محضّرتين عندما تستخدم التجربة هذا الشرط تحديدًا. لذلك نوضح إتاحة المدخلات المحضّرة بصورة منفصلة عن خوارزمية قراءة التعليقات.

أما الرسالة الأصلية وفهارس الموضوعات المقصودة وسجلات قبول المُرسِل، فهي بيانات خاصة بالتقييم. يستنتج الاستخراج المستقل الاختيارات من المشاهدات، بدل أخذها من تلك السجلات. ويختلف عنه مسار المقارنة الذي يستعين بسجل المُرسِل، كما يبين فصل التقييم. وفي نمط التعليقات المتعددة، يتلقى برنامج الاستقبال قائمة مرتبة بمراجع التعليقات الحاملة. أما تجارب الأرشيف، فيبني برنامج التجربة هذه المراجع من الأوقات والمعرّفات الظاهرة، وفق القاعدة الزمنية المتفق عليها.

\subsubsection{من السياق الظاهر إلى الرسالة المستخرجة}
\begin{enumerate}
\item \textbf{تحضير سياق الاستقبال.} ينفذ مسار إعادة البناء تحميل البيانات والبحث في المجال وتوليد الموضوعات باستخدام المدخلات والإعدادات المتفق عليها. وتُستبعد التعليقات الحاملة من السياق السابق للتضمين. وتبني وحدة البتات المشتركة قاموس الضغط. تساعد البصمات الرقمية على مقارنة البنى أثناء التحقق، لكنها لا تحل محل البنى اللازمة للاستخراج.
\item \textbf{قراءة الاختيارين في كل تعليق.} يحدد معرّف التعليق الأب موضع الرد في القائمة المرتبة. ولتحديد الموضوع، يجلب البرنامج الموضوعات القريبة بالبحث الدلالي، ويعيد ترتيبها بحسب الكلمات المشتركة، ثم يطلب من نموذج تحديد الموضوع اختيار فهرس. ويستخدم القائمة وقواعد تحويل الاختيارات نفسها لدى المُرسِل. ويحتفظ مسارا التعليق الواحد والتعليقات المتعددة بالاختلاف بين عروض الحقول الموضح سابقًا.
\item \textbf{جمع البتات وفحصها.} في نمط التعليقات المتعددة، يجمع المُستقبِل المقاطع المرشحة بترتيب التعليقات، ثم يقرأ البادئة والإصدار والعدد والطول. ويفحص \EN{CRC} وأصفار الإكمال، ويشترط تحديد حمولة واحدة قابلة للاستخراج. ويؤدي غياب تعليق أو تعذر حسم التفسير أو إخفاق فحص السلامة إلى نتيجة إخفاق، بدل عرض جزء من الرسالة بوصفه رسالة كاملة.
\item \textbf{استخراج تمثيل الرسالة.} يعكس المُستقبِل تحويل الحمولة المتفق عليه ويفك الضغط وفق التمثيل المستخدم. وتشمل مخرجاته النص المستخرج وحالة النجاح أو الإخفاق، مع معلومات تشخيصية تحدد المرحلة التي وصل إليها. ولا تُقارَن البايتات المستخرجة بالرسالة الأصلية الخاصة إلا بعد انتهاء الاستخراج.
\end{enumerate}

\subsubsection{اختبار المكونات وتجارب استخراج الرسالة كاملة}
تستند صلاحية مكونات المُستقبِل إلى أنواع متعددة من الاختبارات. تختبر اختبارات الوحدات والمكونات تحديد التعليق الأب، وإزالة التعليقات الحاملة من السياق، وتنسيق مراحل إعادة البناء، وتفسير الاختيارات، وجمع التتابع المرتب، وضغط النص وفكه بدقة، ورفض المراجع غير الصالحة والتتابعات المتضررة. وتستخدم بدائل مضبوطة لجلب البيانات واستجابات النموذج لعزل هذه العمليات، فيمكن اختبار معالجة مشاهدة محددة دون استدعاء جديد للشبكة أو النموذج.

ثم تختبر تجارب الرسائل عمليات مترابطة باستخدام تعليقات مولّدة. تختبر تجربة مقتطفات المصادر تحديد الموضوع واستخراج الرسالة كاملة مع مدخلات مشتركة محضّرة. وتضيف تجارب الأرشيف جلب السياق العام وإعادة بناء القائمة بصورة مستقلة في النسخة الحتمية المعتمدة على نصوص المحادثة وحدها. وتفحص اختبارات التطوير، بصورة منفصلة، تطابق تحضير استعلامات البحث والوثائق والموضوعات عند تكراره (القسم~\ref{sec:reconstruction-validation}). ولذلك يستند عمل المُستقبِل إلى اختبارات مكونات وتجارب لاستخراج رسائل كاملة. أما نطاق التكامل المتبقي فهو تجربة واحدة تجمع التحضير الكامل المعتمد على الويب والنموذج مع النشر والاستقبال الفعليين على المنصة؛ ويختلف هذا النطاق عن اختبار كل مكوّن من مكونات المُستقبِل.
